\documentclass[11pt]{article}
\usepackage[margin=1in]{geometry}
\usepackage{amsmath, amssymb, bm}
\usepackage{booktabs}
\usepackage[T1]{fontenc}
\usepackage[utf8]{inputenc}
\usepackage{lmodern}
\usepackage{microtype}
\usepackage{enumitem}
\setlist[itemize]{leftmargin=1.5em}

\title{SAVLG-CBM Loss Terms}
\author{}
\date{\today}

\begin{document}
\maketitle

\section{Model Outputs}
For an input image $x$, the SAVLG model produces:
\begin{itemize}
    \item global concept logits $\bm{c}_{\mathrm{global}} \in \mathbb{R}^K$,
    \item spatial concept maps $\bm{M} \in \mathbb{R}^{K \times H \times W}$,
    \item pooled spatial concept logits $\bm{c}_{\mathrm{spatial}} \in \mathbb{R}^K$,
    \item final concept logits
    \[
        \bm{c}_{\mathrm{final}}
        =
        \bm{c}_{\mathrm{global}}
        +
        \alpha \, \bm{c}_{\mathrm{spatial}},
    \]
    where $\alpha$ is the residual spatial coupling coefficient.
\end{itemize}

The local supervision consists of concept-specific target masks or soft box masks
$\bm{T}_{k} \in [0,1]^{H' \times W'}$ for selected concepts $k$.

\section{Total Objective}
The training objective implemented in \texttt{methods/savlg.py} is
\begin{align}
\mathcal{L}_{\mathrm{total}}
&=
\lambda_{\mathrm{global}} \mathcal{L}_{\mathrm{global}}
+
\lambda_{\mathrm{mask}} \mathcal{L}_{\mathrm{mask}}
+
\lambda_{\mathrm{dice}} \mathcal{L}_{\mathrm{dice}}
+
\lambda_{\mathrm{mil}} \mathcal{L}_{\mathrm{mil}}
\nonumber\\
&\quad
+
\lambda_{\mathrm{outside}} \mathcal{L}_{\mathrm{outside}}
+
\lambda_{\mathrm{coverage}} \mathcal{L}_{\mathrm{coverage}}
+
\lambda_{\mathrm{refine}} \mathcal{L}_{\mathrm{refine}}
+
\lambda_{\mathrm{cons}} \mathcal{L}_{\mathrm{cons}}
+
\lambda_{\mathrm{distill}} \mathcal{L}_{\mathrm{distill}}.
\end{align}

In most successful SAVLG runs, the active subset was
\[
\mathcal{L}_{\mathrm{total}}
=
\lambda_{\mathrm{global}} \mathcal{L}_{\mathrm{global}}
+
\lambda_{\mathrm{mask}} \mathcal{L}_{\mathrm{mask}}
+
\lambda_{\mathrm{outside}} \mathcal{L}_{\mathrm{outside}}
+
\lambda_{\mathrm{coverage}} \mathcal{L}_{\mathrm{coverage}},
\]
optionally with a small Dice term.

\section{Core Loss Terms}

\subsection{Global Concept Loss}
The final concept logits are trained against global concept targets
$\bm{y} \in \{0,1\}^K$ using weighted binary cross-entropy:
\begin{align}
\mathcal{L}_{\mathrm{global}}
=
\frac{
\sum_{k=1}^{K}
w_k^{(\mathrm{global})}
\,
\mathrm{BCEWithLogits}(c_{\mathrm{final},k}, y_k)
}{
\sum_{k=1}^{K} w_k^{(\mathrm{global})}
}.
\end{align}
Positive concepts may receive an upweight through
\texttt{global\_bce\_pos\_weight}.

\paragraph{Purpose.}
This keeps the concept bottleneck predictive at the image level and ensures that
the spatial branch improves the concept vector used by the sparse classifier.

\subsection{Mask Loss}
The mask loss aligns each predicted spatial concept map with its target region.
Let $\bm{m}_k$ be the predicted map logits for concept $k$ after interpolation to
the target mask resolution, and let $\bm{t}_k$ be the corresponding target mask.

SAVLG supports three local loss modes:

\paragraph{(a) BCE mode}
\begin{align}
\mathcal{L}_{\mathrm{mask}}^{\mathrm{bce}}
=
\frac{
\sum_{u}
w_u^{(\mathrm{patch})}
\,
\mathrm{BCEWithLogits}(m_{k,u}, t_{k,u})
}{
\sum_{u} w_u^{(\mathrm{patch})}
}.
\end{align}
This is dense pixel-wise supervision.

\paragraph{(b) Containment mode}
\begin{align}
\mathcal{L}_{\mathrm{mask}}^{\mathrm{contain}}
=
1 -
\frac{
\sum_{u} \sigma(m_{k,u}) \, t_{k,u}
}{
\sum_{u} \sigma(m_{k,u}) + \varepsilon
}.
\end{align}
This only asks that the predicted mass lie inside the target region.

\paragraph{(c) Soft-align mode}
This is the mode used in the strongest runs. The target mask is normalized into a
spatial distribution:
\[
\bm{q}_k = \frac{\max(\bm{t}_k, 0)}{\sum_u \max(t_{k,u},0)}.
\]
The prediction defines a log-softmax distribution:
\[
\log \bm{p}_k = \log \mathrm{softmax}(\bm{m}_k).
\]
Then the loss is
\begin{align}
\mathcal{L}_{\mathrm{mask}}^{\mathrm{soft}}
=
D_{\mathrm{KL}}(\bm{q}_k \,\|\, \bm{p}_k).
\end{align}

\paragraph{Purpose.}
The mask term is the main spatial supervision. In the successful SAVLG runs,
\texttt{soft\_align} worked better than hard per-pixel BCE because it supervises
where probability mass should go without forcing exact mask calibration.

\subsection{Outside-Mass Penalty}
The outside penalty penalizes normalized activation mass that falls outside the
target region. Let $\tilde{\bm{p}}_k = \mathrm{softmax}(\bm{m}_k)$ over spatial
locations. Then
\begin{align}
\mathcal{L}_{\mathrm{outside}}
=
\sum_{u}
\tilde{p}_{k,u}
\,
(1 - \mathrm{clip}(t_{k,u}, 0, 1)).
\end{align}

\paragraph{Purpose.}
This shrinks spillover and encourages the map to stay concentrated inside the
supervised concept region.

\subsection{Coverage Penalty}
The coverage penalty rewards activation over the full target region, not just a
small hotspot. With $\hat{\bm{p}}_k = \sigma(\bm{m}_k)$:
\begin{align}
\mathrm{coverage}_k
&=
\frac{
\sum_{u} \hat{p}_{k,u} \, t_{k,u}
}{
\sum_{u} t_{k,u} + \varepsilon
}, \\
\mathcal{L}_{\mathrm{coverage}}
&=
1 - \mathrm{coverage}_k.
\end{align}

\paragraph{Purpose.}
This addresses the common failure mode where a map peaks at the correct place but
does not cover enough of the ground-truth region to achieve good IoU or box overlap.

\subsection{Dice Loss}
When enabled, Dice loss is
\begin{align}
\mathcal{L}_{\mathrm{dice}}
=
1 -
\frac{
2 \sum_{u} \sigma(m_{k,u}) t_{k,u} + \varepsilon
}{
\sum_{u} \sigma(m_{k,u}) + \sum_{u} t_{k,u} + \varepsilon
}.
\end{align}

\paragraph{Purpose.}
Dice directly optimizes region overlap. In practice, it gave modest gains on the
\texttt{conv5} branch but the biggest localization gains came from changing the
spatial branch itself to \texttt{conv4} or \texttt{conv4+conv5}.

\section{Optional Auxiliary Terms}

\subsection{Local MIL Loss}
If local MIL is enabled, pooled local logits are trained against the global concept
targets using another weighted BCE term:
\begin{align}
\mathcal{L}_{\mathrm{mil}}
=
\mathrm{BCEWithLogits}(\bm{c}_{\mathrm{local\_mil}}, \bm{y}).
\end{align}
This is optional and was usually disabled in the strongest reported runs.

\subsection{Refinement Loss}
\(\mathcal{L}_{\mathrm{refine}}\) is an optional refinement term applied after a
warmup period. It is used to sharpen local supervision but is not part of the core
successful loss recipe described in the main SAVLG runs.

\subsection{Global-Spatial Consistency Loss}
\(\mathcal{L}_{\mathrm{cons}}\) encourages agreement between the global concept
logits and a spatial teacher signal derived from the spatial maps or local MIL logits.
It is optional and activated only if consistency training is enabled.

\subsection{Teacher Distillation Loss}
If a teacher concept model is provided, SAVLG can distill the teacher concept
probabilities using BCE:
\begin{align}
\mathcal{L}_{\mathrm{distill}}
=
\mathrm{BCEWithLogits}(\bm{c}_{\mathrm{final}}, \sigma(\bm{c}_{\mathrm{teacher}})).
\end{align}
This is also optional.

\section{How the Main Successful Configurations Differed}
The main successful settings can be summarized as follows:
\begin{center}
\begin{tabular}{lcccc}
\toprule
Run type & Mask mode & Outside & Coverage & Dice \\
\midrule
conv5 + cov025 & soft-align & 0.25 & 0.25 & 0.00 \\
conv5 + cov050 & soft-align & 0.25 & 0.50 & 0.00 \\
conv5 + dice010 & soft-align & 0.25 & 0.25 & 0.10 \\
conv4 & soft-align & 0.25 & 0.25 & 0.00 \\
conv4+conv5 & soft-align & 0.25 & 0.25 & 0.00 \\
\bottomrule
\end{tabular}
\end{center}

\paragraph{Observed effect.}
Empirically, the largest localization gains came from changing the spatial branch
from \texttt{conv5} to \texttt{conv4} and then to \texttt{conv4+conv5}. Increasing
coverage weight alone did not help, while a small Dice term gave modest gains on
the \texttt{conv5} branch.

\section{Implementation Note}
The formulas in this note are taken directly from the current implementation in
\texttt{methods/savlg.py}, especially \texttt{compute\_spatial\_losses(...)} and
the main training loop where the weighted sum is assembled.

\end{document}
